% File: Volume-I-Foundations/chapters/ch01-intro.tex

\chapter{Introduction and Motivation}
\label{chap:intro}

The Relativistic Scalar--Vector Plenum (\rsvpabbr{}) is a proposal for a
unified field--theoretic framework in which gravity, cosmology, and certain
aspects of information and agency arise from the dynamics of a triplet of
fields
\[
  (\PhiField,\vField,\SField)
\]
living on a spacetime manifold $(M,g)$ or, more generally, on a derived
geometric object that refines $(M,g)$. The scalar field $\PhiField$ encodes a
form of lamphrodynamic potential, $\vField$ represents a flux or flow of
negentropic influence, and $\SField$ is an entropy density that couples to
geometry and matter.

At a heuristic level, \rsvpabbr{} treats spacetime not as a static stage on
which fields live, but as a \emph{plenum} whose local structure is determined
by the coupled evolution of $(\PhiField,\vField,\SField)$. Geometrical
quantities such as curvature, torsion, and metric relations are then derived
from the lamphrodynamic configuration rather than posited \emph{a priori}.

This introductory chapter explains why such a framework naturally leads us to
derived algebraic geometry and shifted symplectic structures, and sets the
stage for the rigorous constructions that follow.

%-----------------------------------------------------------------
\section{From Classical Fields to Derived Geometry}
%-----------------------------------------------------------------

In classical field theory, a field configuration is a section of some
bundle over spacetime $M$, and the space of fields is modeled as an infinite
dimensional manifold or Fréchet space. Functional integrals, symmetries, and
constraints are handled with varying degrees of rigor depending on context.

However, many modern developments --- such as moduli spaces of solutions to
gauge equations, intersection theory on spaces of instantons, and the BV
formalism for gauge theories --- reveal that naive smooth manifolds are
insufficient. One encounters:
\begin{itemize}
  \item Singularities when solutions collide or degenerate.
  \item Higher automorphisms from gauge symmetries.
  \item Obstructions and derived intersections.
\end{itemize}

Derived algebraic geometry was developed precisely to give these moduli spaces
a well--behaved home. Instead of regarding a moduli space as a naive quotient
or set of solutions, one encodes it as a derived stack equipped with additional
structure (cotangent complex, obstruction theory, shifted symplectic form).

\Intuition
Even if the base spacetime $M$ is an ordinary smooth manifold, the \emph{space
of fields} is almost never a manifold in the classical sense. Gauge symmetry,
constraints, and singularities force us to work in a homotopical setting.
Derived stacks and $\infty$--categories are the natural language for this.

In \rsvpabbr{}, lamphrodynamic fields are subject to:
\begin{itemize}
  \item Gauge transformations mixing $\PhiField$ and $\vField$,
  \item Constraints relating $\SField$ to curvature and matter content,
  \item Entropic singularities where $\SField$ becomes non--smooth or diverges.
\end{itemize}
These features suggest that the moduli of lamphrodynamic configurations should
be modeled as a derived stack, which we denote $\Plenum$. A significant portion
of Volume~I is devoted to constructing $\Plenum$, equipping it with a
$(-1)$--shifted symplectic structure, and understanding its tangent and
obstruction theory.

%-----------------------------------------------------------------
\section{The Lamphrodynamic Triplet}
%-----------------------------------------------------------------

We give a preliminary, heuristic description of the fields before formalizing
them in \cref{chap:lamphrodynamic-axioms,chap:plenum-stack}.

\begin{itemize}
  \item $\PhiField$ is a real--valued scalar field on $M$ (or a section of
  some line bundle) representing a lamphrodynamic potential. Regions of high
  $\PhiField$ can be thought of as reservoirs of ``available negentropy''.

  \item $\vField$ is a vector field on $M$, or more generally a section of
  $TM\otimes E$ for some auxiliary bundle $E$, encoding the directed transport
  of lamphrodynamic influence. It plays a role analogous to a current or
  velocity field.

  \item $\SField$ is a scalar entropy density that couples to both geometry and
  matter fields. In contrast to thermodynamic entropy defined on macrostates,
  $\SField$ is intended as a local field that participates dynamically in the
  evolution equations.
\end{itemize}

The basic lamphrodynamic field equations are schematically of the form
\begin{align}
  \partial_t \PhiField &=
    \mathcal{D}_\Phi \Delta \PhiField
    + \mathcal{F}_\Phi(\PhiField,\vField,\SField,g), \label{eq:intro-phi}\\
  \partial_t \SField &=
    \mathcal{D}_S \Delta \SField
    + \mathcal{F}_S(\PhiField,\vField,\SField,g), \label{eq:intro-S}\\
  \partial_t \vField &=
    \mathcal{D}_v \Delta \vField
    + \mathcal{F}_v(\PhiField,\vField,\SField,g), \label{eq:intro-v}
\end{align}
where $\mathcal{D}_\Phi,\mathcal{D}_S,\mathcal{D}_v$ are diffusion constants
and the nonlinear terms $\mathcal{F}_{\Phi,S,v}$ encode lamphrodynamic
interactions and coupling to curvature. Precise versions of these equations
will be derived from an action principle in \cref{chap:aksz-rsvp,chap:variational-hamiltonian}.

\Warning
Equations~\eqref{eq:intro-phi}--\eqref{eq:intro-v} are merely schematic at this
stage. They suggest that \rsvpabbr{} is, at its analytic core, a system of
parabolic or mixed--type PDEs; however, the actual lamphrodynamic equations
will include nonlocal and higher--order terms arising from the BV--AKSZ
construction.

%-----------------------------------------------------------------
\section{Why a $(-1)$--Shifted Symplectic Structure?}
%-----------------------------------------------------------------

In classical Hamiltonian mechanics, the phase space of a system is a symplectic
manifold $(X,\omega)$ and the dynamics are encoded by a Hamiltonian function
$H$ via the equation
\[
  \iota_{X_H}\omega = \dd H.
\]
For gauge theories and field theories, the BV formalism generalizes this
picture to a graded setting: one replaces $(X,\omega)$ by a graded manifold
equipped with an odd symplectic form, and the Hamiltonian $H$ by a master
action $S$ satisfying $\{S,S\}=0$.

In derived geometry, these graded symplectic manifolds are captured by
\emph{shifted symplectic structures} on derived stacks. A $k$--shifted
symplectic structure is, roughly, a closed, nondegenerate $2$--form of
cohomological degree $k$ on the derived stack. The case $k=-1$ is particularly
important for the BV formalism: it corresponds to odd symplectic forms and is
tailored to encode gauge theories.

\begin{definition}
Let $\mathcal{X}$ be a derived Artin stack over a field of characteristic
zero. A \emph{$(-1)$--shifted symplectic structure} on $\mathcal{X}$ is a
closed $2$--form
\[
  \omega \in \mathcal{A}^{2,cl}(\mathcal{X})[-1]
\]
whose underlying pairing on the cotangent complex is nondegenerate in the
derived sense.
\end{definition}

The key insight for \rsvpabbr{} is that the moduli stack $\Plenum$ of
lamphrodynamic configurations should carry such a $(-1)$--shifted symplectic
structure, with the master action functional $S_{\rsvpabbr}$ playing the role
of a Hamiltonian satisfying the classical master equation
\[
  \{S_{\rsvpabbr},S_{\rsvpabbr}\} = 0.
\]
Once this structure is constructed, the AKSZ formalism provides a natural way
to build $S_{\rsvpabbr}$ from geometric data associated to $\Plenum$ and the
spacetime source manifold.

\Intuition
The choice of $k=-1$ is not arbitrary: lamphrodynamic fields are subject to
gauge redundancies and constraints, so their moduli naturally form a BV phase
space rather than a naive symplectic manifold. This forces the symplectic
structure to be of odd degree, hence $k=-1$.

A central technical goal of Volume~I is to:
\begin{itemize}
  \item Construct $\Plenum$ as a derived Artin stack.
  \item Equip $\Plenum$ with a $(-1)$--shifted symplectic structure $\omega$.
  \item Show that $\omega$ can be expressed in terms of $(\PhiField,\vField,\SField)$
  and curvature data, and that it is closed and nondegenerate.
\end{itemize}
These results will be established in \cref{chap:plenum-stack,chap:aksz-rsvp}.

%-----------------------------------------------------------------
\section{Analytic Questions: Well–Posedness and Regularity}
%-----------------------------------------------------------------

The derived geometric formulation of \rsvpabbr{} provides a powerful conceptual
framework, but physics also demands that the associated PDEs be analytically
well--behaved. Concretely, we must address:

\begin{enumerate}
  \item \emph{Well--posedness.} Given initial data $(\PhiField_0,\vField_0,\SField_0)$,
  does there exist a unique solution $(\PhiField,\vField,\SField)$, at least
  for short time, in appropriate Sobolev spaces?

  \item \emph{Energy estimates.} Are there a priori bounds controlling the
  growth of norms of the fields, and do these bounds depend continuously on
  the initial data?

  \item \emph{Regularity and singularities.} Under what conditions can
  singularities form in finite time, and how are such singularities related to
  entropic phenomena in \rsvpabbr{}?

  \item \emph{Long--time behavior.} Do solutions converge to lamphrodynamic
  equilibria or attractors, and can one classify these asymptotic states?
\end{enumerate}

These questions are addressed in Part~III of this volume. We will prove
global well--posedness for linearized lamphrodynamic equations and establish
energy estimates and local well--posedness for certain nonlinear regimes.
Global existence for the fully nonlinear system remains conjectural; we will
state precise open problems and present partial evidence.

%-----------------------------------------------------------------
\section{Organization of the Three Volumes}
%-----------------------------------------------------------------

For orientation, we briefly summarize the roles of the three volumes.

\subsection*{Volume I: Mathematical Foundations}

This volume develops:
\begin{itemize}
  \item The background in derived algebraic geometry and shifted symplectic
  structures needed to formulate \rsvpabbr{} rigorously.
  \item The construction of the lamphrodynamic moduli stack $\Plenum$ and its
  $(-1)$--shifted symplectic structure.
  \item The AKSZ--BV action for lamphrodynamic fields and its classical master
  equation.
  \item Fundamental analytic results on well--posedness, energy estimates, and
  regularity of the lamphrodynamic PDEs.
\end{itemize}

\subsection*{Volume II: Physical Applications}

The second volume applies this foundation to physics:
\begin{itemize}
  \item Derivation of modified Einstein equations with an explicit lamphrodynamic
  stress--energy contribution $\Theta_{ab}$.
  \item Development of a cosmological model in which redshift, structure
  formation, and cosmic microwave background anisotropies are governed by
  lamphrodynamic fields rather than pure metric expansion.
  \item Detailed comparison with observational data (supernovae, BAO, Planck
  CMB measurements) and identification of potential tests that distinguish
  \rsvpabbr{} from $\Lambda$CDM.
\end{itemize}

\subsection*{Volume III: Extensions and Interpretations}

The third volume explores:
\begin{itemize}
  \item The SpherePop calculus as a geometric model of computation, including
  categorical and homotopical formalization and a proof of Turing completeness.
  \item L--systems, semantic fields, and distributed computation (CRDTs) as
  instances of lamphrodynamic or related field theories.
  \item Entropy--weighted spectral geometry and ``cymatic'' resonance modes.
  \item Philosophical and creative dimensions, including the interpretation of
  \rsvpabbr{} within structural realism and process philosophy.
\end{itemize}

%-----------------------------------------------------------------
\section{What is Proven, What is Conjectural}
%-----------------------------------------------------------------

Given the ambitious scope of \rsvpabbr{}, it is important to distinguish
carefully between results that are established in this monograph and those
that remain conjectural.

\begin{itemize}
  \item \emph{Proven in Volume I:} We provide complete proofs for the
  construction of $\Plenum$ as a derived Artin stack, for the existence of a
  $(-1)$--shifted symplectic form on $\Plenum$ under specified conditions, and
  for well--posedness results in the linear and certain nonlinear lamphrodynamic
  regimes.

  \item \emph{Proven in Volume II:} We derive modified Einstein equations from
  the lamphrodynamic action, obtain concrete formulas for cosmological
  observables (including a redshift function $f(\SField,\vField)$), and
  perform quantitative comparisons with existing data. The degree of
  observational agreement is investigated but not treated as a fully
  conclusive fit.

  \item \emph{Proven in Volume III:} We rigorously formalize SpherePop as a
  homotopy--theoretic model of computation, prove that merge operations are
  homotopy colimits, and establish Turing completeness via explicit encodings.

  \item \emph{Conjectural:} Global existence of solutions for the full
  nonlinear lamphrodynamic system; completeness of the cosmological model as a
  full replacement for $\Lambda$CDM; certain speculative connections between
  lamphrodynamic spectral modes and cognition or consciousness.
\end{itemize}

These conjectural aspects are clearly labeled as such, and in each case we
provide reasons for their plausibility and suggest directions for future
research.

%-----------------------------------------------------------------
\section{How to Read this Volume}
%-----------------------------------------------------------------

Different readers may wish to approach the material in different orders.

\begin{itemize}
  \item \textbf{Mathematical physicists} familiar with homological algebra and
  derived geometry may skim Part~I and focus on Parts~II and~III, where the
  lamphrodynamic stack and PDE analysis are developed.

  \item \textbf{Physicists new to derived geometry} are encouraged to read
  \cref{chap:inf-cat,chap:shifted-symplectic} carefully, referring to
  appendices as needed, before tackling the AKSZ construction in
  \cref{chap:aksz-bv,chap:aksz-rsvp}.

  \item \textbf{Analysts} interested in PDE aspects may wish to jump directly
  to Part~III, using the early chapters primarily for motivation and notation.
\end{itemize}

Throughout, we have tried to separate background material, core constructions,
and speculative interpretations. The reader is invited to treat this volume as
a foundation that can be built upon independently of whether they ultimately
accept the full physical picture proposed by \rsvpabbr{}.

\bigskip
In the next chapter, we begin our technical journey by reviewing the basic
language of model categories, $\infty$--categories, and derived stacks, with
an eye toward their eventual application to lamphrodynamic field theory.

